% !TEX root = ../main.tex

\chapter{绪论}
\label{chap:introduction}

\section{研究背景与意义}

关系数据库是组织结构化数据的主要形式之一。在用户行为分析、商品推荐和科学文献管理等场景中，预测所需的信息通常分布于多张相互关联的表中。例如，对用户未来行为的预测可能同时依赖用户属性、历史交互、交互对象及其发生时间。传统机器学习流程通常先指定连接路径、聚合函数和时间窗口，将多表数据整理为单张特征表，再训练预测模型。这一过程需要针对任务设计特征，也预先限定了模型可以利用的关系信息。关系深度学习通过将表中记录表示为节点、主外键关联表示为边，把关系数据库转化为异构图，使跨表信息能够直接参与模型学习\cite{fey2024rdl}。

与一般图数据相比，关系表格数据还具有异构字段和任务时间约束。不同表中的数值、类别和文本字段需要经过相应的编码过程；对带有预测时刻的任务，模型只能使用该时刻之前允许观察的信息。RelBench 提供了公开的关系数据库、预测任务和评价协议，为关系表格学习方法的比较建立了基础\cite{robinson2024relbench}。PyTorch Frame 等工具则提供了处理多模态表格字段的模块化表示\cite{hu2024pytorchframe}。这些工作推动了关系数据库到可训练表示的转换，也使模型表达能力和训练效率成为更突出的研究问题。

从模型表达看，数据库中的主外键可以构成多种具有不同含义的连接路径。例如，学术数据库中的“论文--作者--机构”与“论文--作者--论文”描述了不同的关系。依靠多层图神经网络扩大感受野，需要反复展开和聚合邻域；将不同路径的信息直接混合，也可能使与当前任务有关的语义难以凸显。因此，需要研究如何以适度的模型计算识别、表达并融合跨表关系。

从训练执行看，关系表格数据规模的增长使节点及关联数量持续增加，也使原始字段读取、特征编码和批次构造面临更高的资源开销。采用小批次更新并不意味着只需准备小批次数据。若在训练前加载全部原始表并物化编码特征，主机内存需求随数据库规模增长；若按采样结果逐批读取，则数据库查询、字段转换、文本编码和批次组装可能成为吞吐瓶颈。数据准备速度低于模型消费速度时，训练设备需要等待输入，模型计算的优化难以充分转化为端到端收益。已有图训练系统表明，邻域采样、特征访问、数据传输和训练计算需要作为连续的数据路径加以分析\cite{kaler2022salient,liu2023bgl}。关系表格训练还需从数据库原始行生成异构特征，因此需要研究面向大规模关系表格数据的可扩展训练方法。

本文以关系表格学习为共同研究对象，分别从模型机制和训练框架两个层面开展工作。前者关注跨表语义的有效建模，后者关注数据规模增长时训练过程的可扩展性。模型表达能力与训练执行效率共同影响关系表格学习在大规模数据上的应用。

\section{相关研究现状}

\subsection{关系表格与异构图语义建模}

图表示为多表关系的学习提供了统一入口。GraphSAGE 使用邻域采样和特征聚合支持小批次图训练\cite{hamilton2017graphsage}。关系深度学习进一步利用主外键构建异构图，并结合表格编码器与图模型处理数据库预测任务\cite{fey2024rdl}。在异构图研究中，元路径以节点类型和边类型的序列描述复合关系。HAN 利用注意力机制区分不同元路径的语义贡献\cite{wang2019han}；SeHGNN 将部分邻域聚合移至预处理阶段，以减少训练期间的重复计算\cite{yang2023sehgnn}。

面向关系数据库，RelGNN 根据关系图的结构设计复合消息传递机制\cite{chen2025relgnn}，ReDeLEx 提供了跨多个关系数据库研究不同模型的统一框架\cite{peleska2025redelex}。这些研究分别推进了关系表示、模型结构和统一评价。然而，关系表格学习既要从原始异构字段生成节点特征，又要对不同连接路径承载的语义作出选择。本文第一项工作据此研究一种将初始表示学习与元路径特征聚合分开的模型方法。

\subsection{大规模图训练与数据供给}

大规模图训练的开销不仅来自模型的前向与反向传播。SALIENT 优化邻域采样、特征切片与批次传输，并利用流水线重叠数据准备和设备计算\cite{kaler2022salient}。BGL 关注特征读取、缓存及不同数据准备阶段之间的资源竞争\cite{liu2023bgl}。当数据超出主机内存时，Ginex 利用采样计划组织后续特征访问与缓存\cite{park2022ginex}，MariusGNN 则联合考虑外存布局、数据移动与采样\cite{waleffe2023mariusgnn}。通用机器学习输入流水线的研究同样表明，并行、预取和缓存需要根据实际瓶颈及资源约束进行配置\cite{kuchnik2022plumber}。

上述工作为可扩展训练框架的设计提供了重要依据。直接使用关系数据库训练时，采样结果还需映射回各表中的原始行，再完成异构字段转换和模型输入组装。随着数据库规模扩大，全量特征物化带来更高的内存开销，而逐批读取与编码可能降低训练吞吐。数据库之间的表规模、文本长度、访问重复率和任务时长也具有明显差异。因此，本文第二项工作研究大规模关系表格数据的可扩展训练，以按需数据准备和分阶段执行协调内存占用与端到端训练效率。

\section{研究问题与目标}

本文围绕以下两个问题展开研究。

\begin{enumerate}
  \item \textbf{跨表关系的语义表达。} 同一目标实体可以沿多条主外键路径连接到其他表，不同路径对预测任务的作用并不相同。增加可训练消息传递层数会带来额外的邻域展开和模型计算，却不必然改善对任务相关路径的选择。本文利用元路径表征复合关系，并通过注意力机制融合跨表信息。

  \item \textbf{大规模关系表格数据的可扩展训练。} 随着数据库规模增长，全量物化特征所需的常驻内存持续增加；按需读取虽然降低内存占用，却可能因特征准备速度不足而限制训练吞吐。本文在保持任务划分、时间可见性、采样和特征语义一致的前提下，研究按需特征准备、数据复用与有界流水线，以提高训练过程对数据规模增长的适应能力。
\end{enumerate}

两个问题分别对应模型计算路径与训练执行路径。对第一项工作，本文主要考察预测效果、模型开销和模块作用；对第二项工作，本文主要考察输入正确性、端到端耗时、训练吞吐、整作业内存及不同数据规模下的性能变化。两项工作分别开展实验评估。

\section{主要研究内容}

围绕跨表语义表达问题，本文提出关系表格学习方法 \textbf{MetaRTL}。该方法使用表格编码器与浅层异构图网络进行有限轮次的初始表示学习，再通过非参数传播构造元路径特征，并由 MetaAttn 融合路径语义与全局节点信息。这样的两阶段设计将初始特征学习和任务相关的跨表语义聚合分开。本文通过分类、回归及消融实验分析该方法的建模效果与计算开销。方法细节与实验结果见第\ref{chap:metartl}章。

围绕大规模关系表格数据的可扩展训练问题，本文设计训练框架 \textbf{RelX}。RelX 将可复用的结构与统计准备同在线训练过程区分；在线阶段由采样、特征读取与批次组装、模型训练三个环节构成。不同批次的环节通过有界队列衔接，组装进程按采样计划读取原始行并生成模型输入，结合请求去重和分层复用减少重复工作。RelX 通过控制训练过程的内存占用、减少重复数据准备并重叠不同批次的处理阶段，提高大规模关系表格数据上的训练效率，同时保持既定模型与采样的训练语义。框架设计见第\ref{chap:relx-design}章，正确性、性能与规模扩展实验见第\ref{chap:relx-evaluation}章。

MetaRTL 研究跨表信息的有效表示，RelX 研究大规模关系表格数据的可扩展训练。本文分别报告两项工作的实验结果，并在总结中讨论其适用范围与进一步结合的可能性。

\section{论文组织结构}

全文共六章。第\ref{chap:introduction}章为绪论，介绍研究背景、相关研究、研究问题和主要内容。第\ref{chap:background}章给出关系数据库图表示、异构表格编码、元路径、时序邻域采样与训练流水线等基础概念，并进一步比较相关方法。第\ref{chap:metartl}章介绍 MetaRTL 的模型设计和预测实验。第\ref{chap:relx-design}章介绍 RelX 的总体架构、按需特征准备、缓存与并行执行机制。第\ref{chap:relx-evaluation}章分析 RelX 的正确性、训练效率、内存开销与规模扩展能力。第\ref{chap:conclusion}章总结全文，并讨论研究局限与后续工作。
